\subsection{Cartier descent at an internal radial center}
\label{subsec:internal-cartier-descent}

The finite-root obstruction in
Corollary~\ref{thm:positive-square-picard} depends on the $p$-adic
valuation of the exact logarithmic order. In the singular case of
Theorem~\ref{thm:conductor-equality-structure}, the radial quotient
gives a stronger obstruction: the torsion bundle cannot descend as
an invertible sheaf, and its N\'eron--Severi class is not even a
double. The argument takes place on the smooth resolution of that
quotient. Its polarization contracts the exceptional curve, so the
reducedness step in the Euler-extension argument requires an exact
local multiplicity calculation.

\begin{proposition}\label{prop:internal-cartier-descent-obstruction}
Let $X$ be a smooth projective integral surface over an algebraically
closed field $k$ of characteristic two. Let $Y\subset X$ be a smooth
integral very ample divisor, and let $M\ne\OO_X$ have order two with
$M|_Y\simeq\OO_Y$. Let $2a$ be the exact order along $Y$ of its
logarithmic form, and suppose that
\[
 2g(Y)-2=2aY^2.
\]
Assume that the radial center constructed in the proof of
Theorem~\ref{thm:conductor-equality-structure} lies on $X$, and let
$q:X\to S=X/\mu_2$ be the radial quotient. There is no invertible
sheaf $P$ on $S$ such that $q^*P\simeq M$. Consequently
\[
 M\notin 2\Pic(X),\qquad
 [M]\notin 2\operatorname{NS}(X).
\]
Here $2\Pic(X)$ denotes the subgroup of tensor squares. Equivalently,
the conclusion applies to the case in which the distinguished
normalized torsor has its single $A_1$ singularity.
\end{proposition}

\begin{proof}
\emph{The quotient and its resolution.}
Put $H=\OO_X(Y)$ and $d=Y^2$. We first specify the structural
conclusions that will be used. The proof of
Theorem~\ref{thm:conductor-equality-structure} gives
\[
 K_X\sim(2a-1)Y
\]
and homogeneous coordinates $x_0,\ldots,x_r$ for the embedding by
$H$, with
\[
 v=[1:0:\cdots:0]\in X,
 \qquad Y=X\cap(x_0=0).
\]
The radial action scales $x_0$ and fixes the other homogeneous
coordinates. The divided logarithmic form has no divisorial zero,
and its zero scheme is the single reduced point $v$.

The normal projective quotient $S$ has one singular point
$o=q(v)$, of type $A_1$. Two of the ambient affine coordinates
$u=x_i/x_0$ and $w=x_j/x_0$ are regular parameters at $v$. They
have weight one for the $\mu_2$-action, and the completed invariant
ring is
\begin{equation}\label{eq:internal-cartier-invariant-ring}
 \widehat{\OO}_{S,o}
 =k[[A,B,C]]/(AC-B^2),
 \qquad A=u^2,\quad B=uw,\quad C=w^2.
\end{equation}
Let $b:R=\operatorname{Bl}_vX\to X$ have exceptional curve $F$.
The lifted action has a finite purely inseparable quotient
$q':R\to T$ of degree two, with $T$ smooth and projective. It fits
into
\[
 \rho q'=qb,\qquad \rho:T\longrightarrow S,
\]
where $\rho$ is the minimal resolution, its only exceptional curve
$E$ has square $-2$, and $q'^*E=2F$. On the chart $w=uz$ the
quotient coordinates are $A=u^2,z$, so
\[
 B=Az,\qquad C=Az^2,\qquad E=(A=0).
\]
The saturated lifted foliation on $R$ has line bundle
$\OO_R(b^*Y+F)$. The regular-foliation determinant formula used in
the structure theorem therefore gives
\begin{equation}\label{eq:internal-cartier-canonical-pullback}
 q'^*K_T=K_R-(b^*Y+F)\sim(2a-2)b^*Y.
\end{equation}

\emph{The polarization and Picard pullback.}
The invariant-ring polarization also works with an internal center.
Let $R_X$ be the homogeneous coordinate ring of the given embedding.
Its invariant ring is a graded quotient of
\[
 k[z,x_1,\ldots,x_r],\qquad z=x_0^2,
 \qquad \deg z=2,\quad\deg x_i=1.
\]
Indeed the homogeneous quotient map is equivariant, and its
restriction to each character summand is surjective. The second
Veronese of this weighted polynomial ring is generated in degree
one: each monomial of weighted degree $2n$ is a product of $n$
monomials of weighted degree two. The same holds for its quotient.
Thus $S=\operatorname{Proj}(R_X^{\mu_2})$ has a very ample line
bundle $J=\OO_S(2)$.

There is an isomorphism
\[
 R_X^{\mu_2}/(x_0^2)\simeq R_X/(x_0).
\]
Surjectivity follows from the invariant generators $x_i$ with
$i>0$. If an invariant element is $x_0g$, then $g$ has weight
one, since $x_0$ is a nonzerodivisor. A weight-one polynomial
lifting $g$ is divisible by $x_0$ with invariant quotient. This
proves the kernel assertion. The section $z$ therefore cuts out
a smooth Cartier curve $Y_S\simeq Y$, and
\begin{equation}\label{eq:internal-cartier-polarization}
 J=\OO_S(Y_S),\qquad q^*Y_S=2Y,\qquad Y_S^2=2d.
\end{equation}
The point $o$ is outside $Y_S$. In particular this divisor is
Cartier at $o$, where its section is a unit. Projective normality
of $X$ is not required.

We will use injectivity of Picard pullback for both $q$ and $q'$.
Here is the relevant fixed-point argument, which does not require
the quotient morphism to be flat. Suppose a projective integral
$k$-scheme $W$ with $H^0(W,\OO_W)=k$ has a finite $\mu_p$-quotient
$f:W\to Q$ and a scheme-theoretically fixed $k$-point. If $f^*L$
is trivial, its canonical linearization acts on its one-dimensional
space of global sections by a character. A nonzero section is
nowhere zero. Its value at the fixed point is nonzero, and the
canonical pullback action on that fibre is trivial. The character
is therefore trivial. This invariant section descends, because
\[
 (f_*f^*L)^{\mu_p}
 =L\otimes(f_*\OO_W)^{\mu_p}=L.
\]
Its descent is nowhere zero since $f$ is surjective. Hence $L$
is trivial, proving that $f^*$ on Picard groups is injective.

Every point of $Y$ is scheme-theoretically fixed, and the blowup
does not change a neighborhood of $Y$. The argument applies to
$q$ and $q'$. Put
\[
 D=\rho^*Y_S,\qquad J_T=\OO_T(D)=\rho^*J.
\]
The curve $D$ is smooth, isomorphic to $Y$, and disjoint from $E$.
Equations~\eqref{eq:internal-cartier-canonical-pullback} and
\eqref{eq:internal-cartier-polarization}, together with injectivity
for $q'$, give
\begin{equation}\label{eq:internal-cartier-canonical-class}
 D^2=2d,\qquad q'^*D=2b^*Y,
 \qquad K_T\sim(a-1)D.
\end{equation}
Thus $J_T$ is globally generated, nef and big. We use the pulled-back
space
\[
 V_0=H^0(S,J)\subset H^0(T,J_T)
\]
throughout. It defines the morphism
\[
 f:T\xrightarrow{\rho}S\hookrightarrow\PP(V_0^*).
\]
The only integral curve of intersection zero with
$D$ is $E$: every other curve maps onto a curve on $S$, and its
intersection is positive by the projection formula and ampleness
of $J$.

\emph{The descended logarithmic form.}
Suppose, for a contradiction, that an invertible $P$ on $S$
satisfies $q^*P\simeq M$. Injectivity of $q^*$ gives
\[
 P\ne\OO_S,\qquad P^2\simeq\OO_S,
 \qquad P|_{Y_S}\simeq\OO_{Y_S}.
\]
Its square trivialization gives a global K\"ahler form
$\alpha_S=d\log P$ on $S$, defined by logarithmic differentials
of local units. This definition is valid on the singular surface.
After pullback its square trivialization differs from the chosen
one for $M$ by a scalar, whose differential is zero. Therefore
\begin{equation}\label{eq:internal-cartier-form-pullback}
 q^*\alpha_S=\omega.
\end{equation}
The form $\alpha=\rho^*\alpha_S$ is regular and closed on $T$,
and is nonzero since $q'^*\alpha=b^*\omega$.

Its exact order on $D$ is $a$. To check this, its tangential
restriction is zero because $P|_{Y_S}$ is trivial. The possible
remaining conormal section lies in $\OO_D(-D)$, of negative
degree, and is also zero. Let $n\geq1$ be the exact order. The
leading tangential coefficient is nonzero: a leading section
with zero tangential image would lie in
\[
 H^0(D,\OO_D(-(n+1)D))=0.
\]
At a general closed point of $Y$, the quotient coordinates are
$s=t^2,z=z$. Pullback kills $ds$, while the nonzero leading
tangential coefficient has exact order $2n$. Equation
\eqref{eq:internal-cartier-form-pullback} gives $2n=2a$, hence
$n=a$. This uses completed local rings at a closed point and
makes no perfectness assumption on $k(Y)$.

Let $s_D$ be the canonical section of $J_T$. The form
\[
 \beta=\alpha/s_D^a\in H^0(T,\Omega_T^1(-aD))
\]
is defined everywhere on $T$. Along $D$, divisibility of each
coefficient at its generic point implies divisibility by that
power of its prime local equation in every regular local ring.
Its tangential restriction is nowhere zero on $D$: it is a nonzero
section of a line bundle
of degree
\[
 2g(D)-2-aD^2=0.
\]
Closedness also forces $2\mid a$, by the leading tangential
coefficient calculation. The argument below requires no further
numerical restriction on $a$.

There is no divisorial zero of $\beta$ outside $E$. There is
none on $D$ by the preceding restriction. A zero divisor outside
$D\cup E$ would pull back under the finite map $q'$ to a zero
divisor of $b^*\omega$ away from $b^*Y\cup F$. Such a divisor
would give a divisorial zero of $\omega$ outside $Y$, contrary
to the structural input. The next calculation treats $E$.

\emph{The exceptional multiplicity.}
Choose a local frame of $P$ at $o$. In the completed invariant
ring~\eqref{eq:internal-cartier-invariant-ring}, its square
trivialization is a unit
\begin{equation}\label{eq:internal-cartier-unit-jet}
 U=c+\lambda A+\mu B+\nu C+U_2,
 \qquad c\ne0,\quad U_2\in(A,B,C)^2.
\end{equation}
We have $\mu\ne0$. Indeed, after pulling back to $k[[u,w]]$,
$dA=dC=0$ and
\[
 dB=u\,dw+w\,du.
\]
The radial annihilator and its reduced zero at $v$ make the
first jet of $\omega$ a nonzero scalar multiple of this form.
Indeed the radial derivation has coefficient pair $(u,w)$.
Its annihilating forms are multiples of $u\,dw+w\,du$ by the
Koszul relation, and the reduced zero forces the multiplier to
be a unit.
Terms in $(A,B,C)^2$ have pullback order at least four, and their
differentials cannot contribute to that first jet. Equation
\eqref{eq:internal-cartier-form-pullback} therefore forces
$\mu/c\ne0$.

On the resolution chart $B=Az$, $C=Az^2$, put
\[
 P_2(z)=\lambda+\mu z+\nu z^2.
\]
This polynomial is nonzero, and its derivative is $\mu$. Thus
\begin{equation}\label{eq:internal-cartier-exceptional-form}
 \alpha=
 \frac{(P_2(z)+O(A))\,dA+(A\mu+O(A^2))\,dz}{U}.
\end{equation}
The expansion can be taken in the completed discrete valuation
ring at the generic point of $E$. Substitution sends
$(A,B,C)^2$ into $(A^2)$, so its $z$-derivatives also lie in
$(A^2)$. Completion is faithfully flat, and the displayed orders
are the actual divisorial orders. In particular $\beta$, which
differs from $\alpha$ by a unit near $E$, is nonzero at the
generic point of $E$. It has no divisorial zero anywhere.

Its saturated annihilator is consequently an invertible subsheaf
\begin{equation}\label{eq:internal-cartier-annihilator}
 \mathcal F=\OO_T(aD)\otimes\omega_T^{-1}
 \simeq J_T\longrightarrow T_T.
\end{equation}
This is the Koszul relation between the two local coefficients
of $\beta$, which have no common prime divisor. Near the generic
point of $E$, a generator is a unit multiple of
\begin{equation}\label{eq:internal-cartier-local-field}
 (A\mu+O(A^2))\partial_A
 -(P_2(z)+O(A))\partial_z.
\end{equation}
It is nonzero and tangent to $E$ at its generic point, and its
value on $A$ has order exactly one along $E$.

The composite
\[
 \mathcal F\longrightarrow T_T
 \longrightarrow f^*T_{\PP(V_0^*)}
\]
vanishes along $E$, because $f$ contracts $E$. Its vanishing
order there is exactly one. Indeed $A=u^2$ is the ratio of the
global sections $x_i^2$ and $x_0^2$ of $J$, pulled back to $T$.
Its derivative, computed in
\eqref{eq:internal-cartier-local-field}, has order exactly one.
Every other affine embedding coordinate has derivative divisible
by $A$, since the field is tangent to the contracted curve.
This proves the exact multiplicity assertion for the ambient
tangent map, with its saturated source line bundle.

\emph{An Euler extension on the resolution.}
Pull back the restricted Euler sequence along that composite.
It gives an extension
\[
 0\longrightarrow\OO_T\longrightarrow\mathcal P
 \longrightarrow\mathcal F\longrightarrow0
\]
and a map $\mathcal P\to V_0^*\otimes J_T$.
The middle term is locally free, including along $E$, because
the end terms are line bundles. This does not require the map
to the ambient tangent bundle to be a subbundle at $E$.
Set $\mathcal E=\mathcal P^*\otimes J_T$. By
\eqref{eq:internal-cartier-annihilator} there are maps
\[
 0\longrightarrow\OO_T\longrightarrow\mathcal E
 \xrightarrow{\mathrm{ev}}J_T\longrightarrow0,
 \qquad V_0\otimes\OO_T\longrightarrow\mathcal E,
\]
whose composite is evaluation of sections of $J_T$.
Write $e_s$ for the section induced by $s\in V_0$ and
$e_0=e_{s_D}$. Locally its two components record value and
differentiation in the $\mathcal F$ direction.

The section $e_0$ is nowhere zero. Away from $D$ its value
component is nonzero. On $D$ the normal projection of
$\mathcal F$ is an isomorphism onto $\OO_D(D)$, since the
tangential restriction of $\beta$ is nowhere zero. Its
derivative component is therefore nonzero there. It follows that
\begin{equation}\label{eq:internal-cartier-determinant-extension}
 0\longrightarrow\OO_T\xrightarrow{e_0}\mathcal E
 \xrightarrow{e_0\wedge-}J_T\longrightarrow0
\end{equation}
is exact. The determinant sections $t_s=e_0\wedge e_s$ restrict
on $D$, after one nowhere-zero common factor, to $s|_D$.
On the chart $s_D\ne0$, their orders are those of the derivatives
of $s/s_D$ in the $\mathcal F$ direction.

Choose $s$ with two properties: $s|_D$ has distinct simple zeros,
and $t_s$ has multiplicity exactly one along $E$.
The first is a nonempty open condition, since $J$ embeds the
smooth curve $Y_S$. Explicitly, hyperplanes containing a tangent
line form the image of an incidence variety of smaller dimension
than the dual projective space. The second is a nonempty open
linear condition: all $t_s$ vanish along $E$, and the coordinate
$A$ just exhibited has derivative not vanishing modulo $A^2$.
Both conditions can therefore be satisfied over the infinite
field $k$.

The divisor $C_s=\operatorname{div}(t_s)$ is nonempty and does
not contain $D$. Every component other than $E$ meets $D$, and
the simple zeros of $t_s|_D$ force each such component to have
multiplicity one. The multiplicity of $E$ is one by construction.
A Cartier divisor on the smooth surface $T$ has no embedded
associated components, so $C_s$ is reduced.

It is also connected, although $J_T$ is not ample. If
$C_s=C_1+C_2$ were a disjoint decomposition into nonzero effective
divisors, then
\[
 C_i^2=J_T\cdot C_i\geq0.
\]
A zero square would make $C_i$ orthogonal to the positive-square
class $J_T$. Hodge index would then make $C_i$ numerically zero,
impossible for a nonzero effective divisor because it has positive
intersection with an ample class on $T$. Thus both squares would
be positive; but two positive-square classes cannot be orthogonal
by Hodge index. This contradiction proves connectedness.

The proper reduced connected curve $C_s$ has
$H^0(C_s,\OO_{C_s})=k$. Restricting
\eqref{eq:internal-cartier-determinant-extension} to $C_s$ gives
\[
 e_s|_{C_s}=\gamma e_0|_{C_s}
\]
for a scalar $\gamma\in k$. Dividing $e_s-\gamma e_0$ by
$t_s$ gives a homomorphism $J_T\to\mathcal E$ whose composite
with $e_0\wedge-$ is the identity. Consequently
\[
 \mathcal E\simeq\OO_T\oplus J_T.
\]
For the original evaluation quotient
$\mathrm{ev}:\mathcal E\to J_T$, the component from the first
summand is $s_D$, and the component from $J_T$ is a scalar.
That scalar is nonzero, since otherwise the quotient would
vanish on $D$. Thus the original evaluation quotient splits
as well.

Dualizing and twisting now splits the original Euler pullback
compatibly. The map
\[
 \mathcal F\simeq J_T\longrightarrow f^*T_{\PP(V_0^*)}
\]
lifts to $V_0^*\otimes J_T$. Since $H^0(T,\OO_T)=k$, this
lift is a constant vector $v_1\in V_0^*$. It is nonzero:
away from $E$, the map $f$ is an immersion and the foliation
is generically nonzero.

\emph{The contradiction at the contracted curve.}
The ambient tangent map vanishes along $E$, and $f(E)=o$.
A constant Euler vector has zero projective tangent value at
a point precisely when its projective point is that point.
Therefore $[v_1]=o$. This point is outside the hyperplane
$Y_S$, so the $s_D$-coordinate of $v_1$ is nonzero.

On the affine chart $s_D\ne0$, the induced projective field is
a nonzero scalar multiple of the radial field centered at $o$.
In particular, the affine section ratios $A,B,C$, all
zero at $o$, satisfy
\[
 A\longmapsto c_1A,\qquad
 B\longmapsto c_1B,\qquad
 C\longmapsto c_1C,
 \qquad c_1\ne0.
\]
On the resolution chart $B=Az$, its lift on $A\ne0$ thus sends
$z=B/A$ to zero and $A$ to $c_1A$. These identities hold on the
whole chart by regularity. This lift vanishes along $E$.

On the other hand, the generator of the source line
$\mathcal F$ is, up to a unit, the field in
\eqref{eq:internal-cartier-local-field}. Its $\partial_z$
coefficient is nonzero at the generic point of $E$, since
$P_2(z)\ne0$. The isomorphism $\mathcal F\simeq J_T$ changes
local generators only by units and cannot remove this discrepancy.
The contradiction proves that no invertible $P$ on $S$ pulls
back to $M$.

\emph{Tensor squares and the N\'eron--Severi class.}
Suppose $M\simeq L^{\otimes2}$ on $X$. On invariant affine
charts one has $R^2\subset R^{\mu_2}\subset R$, so relative
Frobenius factors as
\[
 X\xrightarrow{q}S\xrightarrow{g}X^{(2)},
 \qquad gq=F_{X/k}.
\]
The invertible sheaf $P=g^*(L^{(2)})$ is defined also at $o$,
and
\[
 q^*P=F_{X/k}^*L^{(2)}\simeq L^{\otimes2}\simeq M,
\]
contrary to what was just proved. Hence $M$ is not a tensor square.

Finally, if $[M]=2[L_0]$ in $\operatorname{NS}(X)$, then
$M\otimes L_0^{-2}$ is algebraically trivial. The reduced identity
component of the Picard scheme is an abelian variety, and
multiplication by two is surjective on its $k$-points
\cite[Theorems~2.5 and~5.4, Proposition~5.10 and Remark~5.25]{KleimanPicard}
\cite[Tag~03RP, parts~(3)--(4)]{Stacks}.
Choose a square root $L_1$ of that algebraically trivial bundle.
Then $M\simeq(L_0\otimes L_1)^{\otimes2}$, again a contradiction.
This proves the N\'eron--Severi assertion.
\end{proof}

For an internal center, this excludes every double in the
N\'eron--Severi group, independently of $v_2(2a)$. When that
valuation is greater than one, it strengthens the finite-height
obstruction of Corollary~\ref{thm:positive-square-picard}.
The proposition concerns descent as an invertible sheaf on the
singular quotient. A linearization can have different stabilizer
characters on the fixed divisor and at the isolated fixed point;
the proposition does not identify every linearization with such
a descent. Nondivisible N\'eron--Severi classes and positive genus
excess remain outside its conclusion.
